% TOP_PROBLEM: {"problemId": "problem.p-versus-pspace", "canonicalTitle": "P versus PSPACE", "releaseRank": 14, "primaryDomain": "theoretical_computer_science", "primaryDomainLabel": "Theoretical computer science", "displayStatus": "open", "releaseStatus": "open"}
% !TeX program = lualatex
% ATTEMPT_STATUS: unresolved
\documentclass[11pt]{article}
\usepackage[margin=25mm]{geometry}
\usepackage{fontspec}
\setmainfont{Latin Modern Roman}
\usepackage{amsmath,amssymb,mathtools}
\usepackage{unicode-math}
\setmathfont{Latin Modern Math}
\usepackage{xurl,hyperref}
\hypersetup{colorlinks=true,urlcolor=blue}
\setcounter{secnumdepth}{0}
\setlength{\parindent}{0pt}
\setlength{\parskip}{5pt}
\setlength{\emergencystretch}{3em}
\begin{document}
\section*{\texorpdfstring{P versus PSPACE}{P versus PSPACE}}\label{14}
\subsection{Definitions and mathematical statement}
For an input of length $n$, time counts computation steps and space counts work-tape cells. Define
\[
 \mathsf P=\bigcup_{k\ge1}\mathsf{DTIME}(n^k),\qquad
 \mathsf{PSPACE}=\bigcup_{k\ge1}\mathsf{DSPACE}(n^k).
\]
Machines must halt and decide their language correctly on every input. Determine whether these classes coincide. An algorithm with polynomial space and exponential time does not establish membership in $\mathsf P$.
\par\smallskip\textit{Formulation and concept references:} \href{https://art.torvergata.it/bitstream/2108/32769/60437/ComputabilityDecidabilityComplexity.pdf}{[S1]}.

\subsection{Short English statement}
Can every decision problem solvable with polynomially much memory also be solved in polynomially much time?
\subsection{Research attempt}
\paragraph{Scope and literature check.}
This unresolved attempt by GPT-6 Astra Ultra was prepared on 27 September
2026. Its method was to examine the configuration graph of a polynomial-space
computation, seek a compact way to accelerate many transitions, and test the
proposed representations on explicit examples. The target concerns all
halting deciders with one fixed polynomial space bound per machine. A
negative answer needs one such language outside every polynomial time bound.

Williams [S2], Theorem 1.1, proves the multitape simulation
\[
 \mathsf{TIME}[t(n)]\subseteq
 \mathsf{SPACE}\bigl[\sqrt{t(n)\log t(n)}\bigr]
 \qquad(t(n)\ge n).
\]
Corollary 1.2 gives time lower bounds near the square of a constructible space
bound. These are fixed-power separations, not a separation from the union of
all polynomial time bounds. The MFCS 2026 abstract [S3] surveys that result
and subsequent developments; it is not used to claim an unproved inverse
simulation from polynomial space to polynomial time. The configuration-graph
background is also discussed in [E1], Chapter 4. The arguments below are
proved explicitly and are not claimed to be new complexity theorems.

\paragraph{Attack plan: compress a deterministic trajectory.}
A deterministic computation using $s$ configuration bits has at most $2^s$
configurations. A halting trajectory cannot repeat a nonhalting configuration:
once repeated, determinism would force a loop. Thus exponential simulation
is available with polynomial memory, but the central issue is whether the
trajectory can be skipped in large blocks.

Two candidate block representations are a circuit for the iterated successor
and an algebraic map over $\mathbb F_2$. The affine case admits a complete
polynomial-time algorithm for an absorbing target. General Boolean circuits
already encompass the full original difficulty; low-degree coordinate
polynomials need not stay small when expanded under composition.

\paragraph{Lemma 14.1: the succinct functional-graph target.}
Consider the following language of encoded triples $(F,a,t)$. The input
contains a Boolean circuit for a total map
$F:\{0,1\}^s\to\{0,1\}^s$, together with $a,t\in\{0,1\}^s$.
Accept exactly when $F(t)=t$ and $F^j(a)=t$ for some integer $j\ge0$.
This language is PSPACE-complete under polynomial-time many-one reductions.

\emph{Membership.} Check $F(t)=t$. Starting at $a$, evaluate successive
states and test whether they equal $t$. If the first $2^s$ states, indexed
$0$ through $2^s-1$, omit $t$, some state repeats before a first possible
later visit, and determinism rules that visit out. A binary counter with
$s+1$ bits suffices to implement this bound. Circuit evaluation and the
current state require space polynomial in the input encoding length.
Therefore the decision procedure is in PSPACE; its time need not be
polynomial.

\emph{Hardness.} Fix a language $L$ in PSPACE and a halting deterministic
machine $M$ deciding it in at most $p(n)$ work cells. On a fixed input $x$,
encode a configuration by its work-tape contents, head positions, and state,
using $s=\operatorname{poly}(|x|)$ bits. The fixed input $x$ is hardwired
into a polynomial-size circuit computing one transition. Validity of the
configuration encoding and head positions is also checked by polynomial-size
circuits. Reserve two additional tagged configurations for acceptance and
rejection. Map an accepting machine configuration to the acceptance tag,
a rejecting one to the rejection tag, and both tags to themselves. Invalid
encodings can map to the rejection tag. These tests and updates are local
or polynomial-size comparisons, so the resulting total circuit is
constructible in polynomial time in $|x|$. Let $a$ encode the initial
configuration and $t$ be the acceptance tag. The resulting trajectory reaches
$t$ precisely when $M$ accepts $x$, and $t$ is absorbing. This proves the
reduction.\hfill$\square$

Thus a general polynomial-time algorithm for this particular succinct
functional-graph problem would prove $\mathsf P=\mathsf{PSPACE}$. A long
trajectory alone would not prove the opposite: the decision problem allows
an algorithm to determine the destination without visiting every state.
For example, a binary counter takes exponentially many increments to reach
a specified high count, but its order and destination are easily understood
arithmetically.

\paragraph{Lemma 14.2: a complete affine special case.}
Suppose $F(v)=Av+b$ over $\mathbb F_2$, where
$A\in\mathbb F_2^{s\times s}$ and $b\in\mathbb F_2^s$ are explicitly given,
and suppose $F(t)=t$. Then
\begin{equation}\label{a14:affine}
 \exists j\ge0:\ F^j(a)=t
                  \quad\Longleftrightarrow\quad A^s(a-t)=0.
\end{equation}
The right side is computed in $O(s^3)$ field operations and $O(s^2)$
working bits by at most $s$ matrix-vector multiplications.

\emph{Proof.} Since $At+b=t$,
$F(v)-t=A(v-t)$, and induction gives
$F^j(a)-t=A^j(a-t)$. Write $K_j=\ker A^j$, with $K_0=\{0\}$.
The subspaces satisfy $K_j\subseteq K_{j+1}$. If
$K_j=K_{j+1}$, then $A^{j+2}v=0$ implies
$Av\in K_{j+1}=K_j$, whence $v\in K_{j+1}$. Thus
$K_{j+2}=K_{j+1}$ and the chain stays constant thereafter. There can be
at most $s$ strict dimension increases, so it has stabilized by $K_s$.
Consequently $A^j(a-t)=0$ for some $j$ exactly when $a-t\in K_s$.
The computation bound follows from $s$ products, each costing $O(s^2)$
field operations.\hfill$\square$

The absorbing-target hypothesis is essential to this proof. An arbitrary
orbit-membership question for affine maps is not being silently identified
with kernel membership. Moreover the result uses explicit coefficients;
no algorithm for recognizing or discovering a hidden affine conjugacy of an
arbitrary succinct graph has been supplied.

A useful interpretation is that the part of an affine orbit that can die
at a fixed point is governed by an ascending chain of kernels. Its dimension
is the potential function limiting the relevant evolution. This succeeds
even though the same affine map may have long cycles elsewhere. A potential
function for general Boolean transitions with an analogous polynomial bound
would be a substantive missing ingredient.

\paragraph{Testing polynomial composition as a generalization.}
Every Boolean function has a unique multilinear polynomial over
$\mathbb F_2$. It is tempting to store each successor coordinate in this
form and square the transition by substitution, reducing $z^2$ to $z$.
The following family shows that an expanded monomial representation can
already become exponential after logarithmically many simple updates.

Let $k$ be a power of two. Use $2k$ fixed leaf coordinates
$x_1,y_1,\ldots,x_k,y_k$, a coordinate $u_i$ for each pair, and one coordinate
for every internal node of a balanced binary tree with leaves $u_i$.
Define a synchronous map $G$ as follows: the $x_i,y_i$ coordinates stay fixed,
$u_i$ is replaced by $x_i+y_i$, and each tree-node coordinate is replaced by
the product of its two child coordinates. There are $4k-1$ coordinates and
each update polynomial has degree at most two. The entire map has a
Boolean circuit of size $O(k)$.

After $1+\log_2 k$ iterations, the root coordinate is
\begin{equation}\label{a14:expansion}
                   \prod_{i=1}^k(x_i+y_i)
       =\sum_{\varepsilon\in\{0,1\}^k}
          \prod_{i=1}^k
          \begin{cases}x_i,&\varepsilon_i=0,\\y_i,&\varepsilon_i=1.
          \end{cases}
\end{equation}
This follows by induction on the level of the tree: after the first update
the leaves of that tree contain the desired pair sums, and each successive
update propagates their products one level. Initial values of the auxiliary
coordinates no longer affect the root at the stated time. The $2^k$
monomials in the expansion are squarefree and distinct, since each selects
one variable from each disjoint pair. None cancel over $\mathbb F_2$.
Uniqueness of the multilinear polynomial therefore forces $2^k$ stored
monomials in an explicit expansion.

This falsifies the proposed universal polynomial bound for that representation.
It does \emph{not} lower-bound the size of circuits for the iterates: the
factored expression on the left has a circuit with $O(k)$ gates. Keeping
factorization repairs this example. The remaining question is whether such
compact factorizations or other summaries can always be found and manipulated
for arbitrary transitions. Simply writing the formal composition
$F^{2^{r+1}}=F^{2^r}\circ F^{2^r}$ supplies two copies of the prior circuit,
so the direct construction doubles its size at each stage. It gives no
uniform polynomial bound after $s$ stages.

\paragraph{Verification and boundary cases.}
The companion script exhaustively checks \eqref{a14:affine} for every binary
matrix of dimensions $1,2,3$, every translation $b$, every fixed point $t$,
and every starting point $a$, comparing the kernel test with direct orbit
exploration until repetition. It also expands \eqref{a14:expansion} for
$k=1,2,4,8$, verifies $2^k$ distinct monomials, and compares its truth values
with the factored expression through $k=4$. The tests are finite checks of
these lemmas, not checks of the open separation.

If $a=t$, the algorithm accepts immediately and $A^s(a-t)=0$. A permutation
matrix has zero kernel, so only $a=t$ can reach the absorbing target; a
nilpotent matrix reaches it from every translated starting point. These
opposite cases agree with the lemma. The circuit reduction has only one
successor per state, so no nondeterminism is concealed in its hardness.

\paragraph{Remaining gap and outcome.}
\textbf{UNRESOLVED.} The proved special-case algorithm uses a kernel chain
whose dimension bounds all progress toward a fixed point. The attempted
expanded-polynomial extension fails on an explicit linear-size family.
A positive solution still needs polynomial-time reachability for arbitrary
succinct deterministic transitions, or a different equivalent algorithm.
A negative solution needs a lower bound against all polynomial-time deciders,
not merely against simulation or monomial expansion. Concrete continuations
would seek a richer invariant preserved under composition, bound the total
size and cost of its updates, and test it on nonlinear circuits encoding
polynomial-space machines. None of those general bounds has been established.

\subsection{Sources}
\begin{itemize}
\item[S1]Walter S. Brainerd and Lawrence H. Landweber et al., Computability, Decidability, Complexity, third edition.
\url{https://art.torvergata.it/bitstream/2108/32769/60437/ComputabilityDecidabilityComplexity.pdf}.
\textit{Formulation source.}
\item[S2]Simulating Time With Square-Root Space — Ryan Williams, February 24, 2025.
\url{https://eccc.weizmann.ac.il/report/2025/017/download}.
\textit{Further reference; not a proof endorsement.}
\item[S3]Some Recent Developments in Space Complexity — R. Ryan Williams, MFCS 2026.
\url{https://drops.dagstuhl.de/entities/document/10.4230/LIPIcs.MFCS.2026.3}.
\textit{Further reference; not a proof endorsement.}
\item[S4]Some conditions implying if P=NP then P=PSPACE — Ismael Rodriguez (February 10, 2026).
\url{https://arxiv.org/abs/2602.10073}.
\textit{Further reference; not a proof endorsement.}
\item[E1]Sanjeev Arora and Boaz Barak,
\emph{Computational Complexity: A Modern Approach}, web draft,
8 January 2007, Chapter 4, space complexity and configuration graphs.
\url{https://theory.cs.princeton.edu/complexity/book.pdf}.
\textit{Background definitions; consulted 27 September 2026. The restricted
functional-graph reduction and affine algorithm used here are proved in
this attempt.}
\end{itemize}
\end{document}
